\section{Proofs for Section~\ref{sec:growth}}\label{app:growth}

\begin{proof}[Proof of Lemma~\ref{lem:graded}]
(a) Consider the side towards $u_i'$; the other is symmetric. The grid interval
between the nodes at distances $t_k<t_{k+1}$ has length at most $\sigma(t_k)$.
For $i\in I_C$ this is $h+\theta t_k$. For $i\in I_Z$ and a length of at least
two, $\sigma(t_k)\ge2$, so $\sigma(t_k)=\lfloor h+\theta t_k\rfloor\le h+\theta t_k$;
intervals of length one are ignored in $w_i$. The node at distance $t_k$ has
adjacent intervals starting at distances $t_{k-1}$ and $t_k$, both at most
$t_k=\abs{v-c_i}$; the intervals adjacent to the center have length at most $h$.

(b) Let $t_1<\dots<t_m=R$ be the nodes on that side. Every step except the last
is unclipped, and an unclipped step from $t$ has length
$\sigma(t)\ge(H+\theta t)/3$: for $i\in I_C$ trivially; for $i\in I_Z$ with $h\ge1$
because $\lfloor y\rfloor\ge y/2$ for $y\ge1$; for $i\in I_Z$ with $h<1$, either
$\theta t<2$ and $\sigma\ge1>(1+\theta t)/3$, or $\theta t\ge2$ and
$\lfloor h+\theta t\rfloor\ge\theta t/2\ge(1+\theta t)/3$. Hence
$t_{k+1}+H/\theta\ge(1+\theta/3)(t_k+H/\theta)$ for $k\le m-2$, and
$(1+\theta/3)^{m-1}\le(t_{m-1}+H/\theta)\theta/H<1+\theta R/H$. So
$m-1<\ln(1+\theta R/H)/\ln(1+\theta/3)$, and
$\ln(1+x)\ge\frac{23}{24}x$ for $0\le x=\theta/3\le\frac1{12}$ gives the bound.

(c) Each side has length $R\le h$, and $\sigma(0)\ge R$ (for $i\in I_Z$, $R$ is
an integer and $\lfloor h\rfloor\ge R$), so the first step reaches the endpoint.
\end{proof}

\paragraph{Bit lengths in Theorem~\ref{thm:approx}.}
Let $\Gamma_X$ be the product of the denominators of the box endpoints and
$\Gamma_F$ that of the coefficients, after substituting fixed coordinates;
both have $O(I)$ bits, and $\abs{e_i},\abs E=O(I)$. Fix a trial $\mu$ and let
$\alpha=J+\max_i\abs{e_i}+\abs E+\mu K_\mu$. Every node of every grid of the
trial has a denominator dividing $\Gamma_X2^\alpha$. Indeed, the initial center
and the original endpoints have denominators dividing $\Gamma_X$, and a new node
is a center coordinate plus or minus
\[
h_{ij}\,\frac{(1+\theta)^k-1}{\theta}
=2^{E-j-e_i}\,\frac{(2^\mu+1)^k-2^{\mu k}}{2^{\mu(k-1)}},\qquad k\le K_\mu+1,
\]
or an integer step for $i\in I_Z$, or a clipped endpoint of the current box;
centers and box endpoints are nodes of earlier stages of the same trial or
original data. So denominators do not grow with the number of stages. Nodes
lie in $\bar X$, so their numerators have $O(I+\alpha)$ bits. Every factor
value, correction $L_iw^2/8$, table entry and message has a denominator dividing
$8\Gamma_F\Gamma_X^24^\alpha$ and is a sum of at most $\abs{\mathcal A}+n$ terms
of the form coefficient times a product of at most two nodes (a message is the
value of such a sum at a minimizing assignment). Over this common denominator
every number is an integer of $b=O(I+\alpha)$ bits, the dynamic program uses
only integer additions and comparisons, and a factor evaluation costs $O(b^2)$
bit operations per monomial. Since $J+1\le c(I+q+1)$ and
$\lceil\log_2(n_P+2)\rceil\le I+1$, we have $\mu K_\mu\le20\mu2^\mu(I+1)$ and
$b\le c(1+\mu2^\mu)(I+q+1)$. With $2^{\mu^*}\le6\sqrt{\bar\kappa}$ and
$\mu^*\le3(1+\log_2\bar\kappa)$, the table work of Theorem~\ref{thm:approx} times
$b^2$ gives the stated $f(p,\bar\kappa)(I+q+1)^5$. The certificate stores the
grids of one trial and is checked by the same dynamic programs. Without growth
the schedule still halts, and every number it forms is bounded as above for
its trial.

\begin{proof}[Proof of the claims in Example~\ref{ex:family}]
(a) Let $d=(1-u,1-v)\in[0,1]^2$. At $(1,1)$ the gradient of
$u^2+v^2-4uv+\frac{u+v}4$ is $(-\frac74,-\frac74)$ and its quadratic part is
$d_1^2+d_2^2-4d_1d_2\ge-(d_1^2+d_2^2)$, so, using $d_i\ge d_i^2$,
\[
u^2+v^2-4uv+\tfrac{u+v}4+\tfrac32\ge\tfrac74(d_1+d_2)-(d_1^2+d_2^2)
\ge\tfrac34(d_1^2+d_2^2).
\]
With $(\rho-\frac12)^2\le2(\rho-\frac u2)^2+\frac12d_1^2$ this gives
$\phi-\phi^*\ge\frac34(d_1^2+d_2^2)+(\rho-\frac u2)^2\ge\frac12\bigl(d_1^2+d_2^2
+(\rho-\frac12)^2\bigr)$. The coupling is nonnegative and vanishes at $x^*$;
summing over blocks gives (a). (b) The diagonal Hessian entries are
$\frac52+\frac{2\deg(b)}{16\Delta}$ for $u_b$ and $2$ for $v_b,\rho_b$.
(c) At such a point $\bar x$ the coupling gradient in $u_b$ has absolute value at
most $\frac18$. Hence $\partial_{u_b}F\ge\frac18$ and $\partial_{v_b}F=\frac14$ at
$u_b=v_b=0$, $\partial_{u_b}F\le-\frac74+\frac18$ and $\partial_{v_b}F=-\frac74$ at
$u_b=v_b=1$, and $\partial_{\rho_b}F=0$. For feasible $x$ let $d=x-\bar x$ and
$\delta=\sum_b(\abs{d_{u_b}}+\abs{d_{v_b}})$. The first-order term is at least
$\delta/8$; the second-order term is at least
$\sum_b(d_{\rho_b}-\frac12d_{u_b})^2-\delta^2$. Thus
$F(x)-F(\bar x)\ge\delta/8-\delta^2+\sum_b(d_{\rho_b}-\frac12d_{u_b})^2>0$ for
$0<\delta<\frac18$, and $F(x)-F(\bar x)=\sum_bd_{\rho_b}^2>0$ for $\delta=0\ne d$.
(d) Attach to a bag containing $u_b$ a new bag $\{u_b,v_b,\rho_b\}$.
\end{proof}

\begin{proof}[Proof of Proposition~\ref{prop:sharp}]

The Hessian is block diagonal with blocks
$\bigl(\begin{smallmatrix}\Lambda&\sigma\\\sigma&\Lambda\end{smallmatrix}\bigr)$,
whose eigenvalues are $\Lambda\pm\sigma>0$. Hence $F(x)\ge\frac{\Lambda-\sigma}2
\norm x^2$ with equality along $u_k=-v_k$, so $x^*=0$, $\OPT=0$ and
$g=(\Lambda-\sigma)/2$. Each diagonal entry is $\Lambda$. Since
$\kappa=2\Lambda/(\Lambda-\sigma)$ increases continuously from $2$ to
$\infty$ as $\sigma$ goes from $0$ to $\Lambda$, every $\kappa\ge2$ occurs.

The corrected objective separates over pairs,
$Q=\sum_kq_k(u_k,v_k)$ with
$q_k(u,v)=\frac\Lambda2(u^2+v^2)+\sigma uv-\frac\Lambda8(w(u)^2+w(v)^2)$,
where $w$ denotes the largest adjacent interval length in the relevant
grid. Therefore the min-marginal of $u_1$ is
$m_{u_1}(a)=\min_vq_1(a,v)+\sum_{k\ge2}\min q_k$, the minima taken over the
grids. By hypothesis $w(0)\ge h$ in every coordinate, so
$\min q_k\le q_k(0,0)\le-\Lambda h^2/4$ for each $k\ge2$, and
$\sum_{k\ge2}\min q_k\le-(m-1)\Lambda h^2/4=-(n-2)\Lambda h^2/8$.

Fix a node $a$ with $|a|\le\rho$ and put $t^*=-\sigma a/\Lambda$; then
$|t^*|\le|a|\le\rho$, so $t^*$ lies in the box of $v_1$ and in some grid
interval of $G_{v_1}$ of length $\lambda$, say. Let $v_a$ be the endpoint of
that interval nearest to $t^*$. Then $|v_a-t^*|\le\lambda/2$ and
$w(v_a)\ge\lambda$. Completing the square,
$\frac\Lambda2v^2+\sigma av=\frac\Lambda2(v-t^*)^2-\frac{\sigma^2a^2}{2\Lambda}$,
so
\[
q_1(a,v_a)\le\frac{\Lambda^2-\sigma^2}{2\Lambda}a^2
+\frac\Lambda2\Bigl(\frac\lambda2\Bigr)^2-\frac\Lambda8\lambda^2
=\frac{\Lambda^2-\sigma^2}{2\Lambda}a^2 ,
\]
where we also dropped the nonpositive term $-\frac\Lambda8w(a)^2$. Since
$\Lambda-\sigma=2g$ and $\Lambda+\sigma=2(\Lambda-g)$,
$(\Lambda^2-\sigma^2)/(2\Lambda)=2g(\Lambda-g)/\Lambda$. Thus
\[
m_{u_1}(a)\le\frac{2g(\Lambda-g)}{\Lambda}a^2-\frac{(n-2)\Lambda h^2}8,
\]
which is at most $0\le U$ when
$a^2\le(n-2)\Lambda^2h^2/(16g(\Lambda-g))=(n-2)h^2\kappa^2/(16(\kappa-1))$.
Because $\kappa^2/(\kappa-1)\ge\kappa$, the hypothesis
$a^2\le(n-2)\kappa h^2/16$ suffices. The coordinates $v_k$ and the other
pairs are symmetric. An interval with an endpoint $a$ satisfying
$m_i(a)\le U$ passes the filtering test.
\end{proof}
\begin{proof}[Proof of Corollary~\ref{cor:uniformgrid}]

Uniform rule, center. Suppose a stage is centered at $0$ with box
$B^\circ\ni0$. Its grid in coordinate $i$ consists of the multiples of $h_j$
in $B^\circ_i$ and the endpoints of $B^\circ_i$. All grid intervals have length
at most $h_j$, and the longest one adjacent to $0$ is at least as long as any
other interval, because only the outermost intervals are shortened by the
endpoints. Hence $w_i(v)\le w_i(0)$ for all nodes, so for every grid point
$x\ne0$,
$Q(x)-Q(0)\ge F(x)-F(0)\ge\frac{\Lambda-\sigma}2\norm x^2>0$. The corrected
minimizer is $0$, and so is the next center; by induction every stage is
centered at $0$.

Uniform rule, retained box. Put $\rho_j=\min\{1,2(r+1)h_j\}$; we show by
induction that the box $B_j$ of stage $j$ contains $[-\rho_j,\rho_j]^n$.
This holds for $j=0$ ($B_0=[-1,1]^n$). At stage $0$ ($h_0=2$) the grid is
$\{-1,0,1\}$ in every coordinate, with $w_i(0)=1$; by
Proposition~\ref{prop:sharp} with $h=1$ the node $0$ passes, so both
intervals are retained and $B_1=[-1,1]^n\supseteq[-\rho_1,\rho_1]^n$. Let
$j\ge1$. Then $h_j\le1$, hence $h_j\le\rho_j$, so $\pm h_j$ are nodes and
$w_i(0)=h_j$. Proposition~\ref{prop:sharp} with $h=h_j$ and
$\rho=\rho_j$ retains every interval with an endpoint $kh_j$, where
$|k|\le r$ and $|k|h_j\le\rho_j$. If $(r+1)h_j\le1$, then $(r+1)h_j\le\rho_j$
and the retained intervals $[kh_j,(k+1)h_j]$, $-r-1\le k\le r$, give
$B_{j+1}\supseteq[-(r+1)h_j,(r+1)h_j]^n=[-\rho_{j+1},\rho_{j+1}]^n$. If
$(r+1)h_j>1$, then $\rho_j=1$, $B_j=[-1,1]^n$, and the retained intervals
reach the box boundary (the outermost ones possibly shortened), so
$B_{j+1}=[-1,1]^n\supseteq[-\rho_{j+1},\rho_{j+1}]^n$. This completes the
induction. When $(r+1)h_j\le1$, the grid of stage $j+1$ has spacing $h_j/2$
and contains every multiple of $h_j/2$ in $[-(r+1)h_j,(r+1)h_j]$, that is,
$4(r+1)+1$ nodes; and $r+1\ge\sqrt{(n-2)\kappa/16}$ gives
$4r+5\ge\sqrt{(n-2)\kappa}+1$.

Graded rule. Apply Proposition~\ref{prop:sharp} with $h=h_j$ and
$\rho=R$. Every point $t\in[-R,R]$ lies in a grid interval $[a,b]$; if
$t\ge0$ then $0\le a\le R$ (the node $0$ lies in the grid), and if $t<0$ then
$-R\le b\le0$, so the interval has an endpoint of absolute value at most
$R$ and is retained. Hence the retained box contains $[-R,R]$ in every
coordinate. The next grid starts from its center $c'$ inside this box and
uses unclipped steps $h'+\theta t$ with $h'=h_j/2$, so after $k$ steps the
distance from $c'$ is at most $h'((1+\theta)^k-1)/\theta$. One side of $c'$
has length at least $R$ within the box, and must be covered; hence
$k\ge\log(1+\theta R/h')/\log(1+\theta)$ nodes lie on that side, besides
$c'$.
\end{proof}
